\documentclass[14pt,a4paper]{extarticle}

% Кодировка и язык (XeLaTeX)
\usepackage{fontspec}
\setmainfont{Times New Roman}
\setsansfont{Helvetica}
% monospace: use default CM font for Cyrillic compatibility
\usepackage{polyglossia}
\setdefaultlanguage{russian}
\setotherlanguage{english}
% Suppress monospace Cyrillic warnings
\newfontfamily\russianfonttt{Times New Roman}

% Поля по ГОСТу
\usepackage[left=30mm,right=15mm,top=20mm,bottom=20mm]{geometry}

% Межстрочный интервал 1.5
\usepackage{setspace}
\onehalfspacing

% Абзацный отступ 1.25 см
\usepackage{indentfirst}
\setlength{\parindent}{1.25cm}

% Глобальная защита от переполнения строк
\tolerance=9000
\emergencystretch=3em
\hbadness=10000
\setlength{\overfullrule}{0pt}

% Изображения
\usepackage{graphicx}
\usepackage{float}

% Таблицы
\usepackage{array}
\usepackage{tabularx}
\usepackage{longtable}

% Типы колонок для таблиц
\newcolumntype{L}[1]{>{\raggedright\arraybackslash}p{#1}}
\newcolumntype{C}[1]{>{\centering\arraybackslash}p{#1}}

% Подписи
\usepackage{caption}
\captionsetup[figure]{name=Рисунок,labelsep=endash}
\captionsetup[table]{name=Таблица,labelsep=endash}

% Нумерация страниц
\usepackage{fancyhdr}
\pagestyle{fancy}
\fancyhf{}
\fancyfoot[C]{\thepage}
\renewcommand{\headrulewidth}{0pt}

% Гиперссылки
\usepackage[hidelinks]{hyperref}
\usepackage{url}

% Диаграммы
\usepackage{tikz}
\usetikzlibrary{arrows.meta, positioning, shapes.geometric, fit, calc}

\begin{document}

% ========== ТИТУЛЬНЫЙ ЛИСТ ==========
\begin{titlepage}
\thispagestyle{empty}
\begin{center}

{\bfseries Министерство науки и высшего образования Российской Федерации}

\vspace{2mm}

{\small\bfseries ФЕДЕРАЛЬНОЕ ГОСУДАРСТВЕННОЕ АВТОНОМНОЕ ОБРАЗОВАТЕЛЬНОЕ\\
УЧРЕЖДЕНИЕ ВЫСШЕГО ОБРАЗОВАНИЯ}

\vspace{2mm}

{\bfseries <<Национальный исследовательский университет ИТМО>>\\
(Университет ИТМО)}

\vspace{8mm}

\begin{flushleft}
\hspace{1.25cm}\begin{minipage}[t]{\dimexpr\textwidth-1.25cm\relax}
\textbf{Факультет:} Инфокоммуникационных технологий

\textbf{Образовательная программа:} Инфокоммуникационные системы

\textbf{Направление подготовки (специальность):} Прикладное программирование в инфокоммуникационных системах
\end{minipage}
\end{flushleft}

\vfill

{\Large\bfseries ИТОГОВЫЙ ОТЧЁТ}

\vspace{4mm}

по дисциплине <<Создание клиент-серверных приложений>>

\vspace{4mm}

Тема: \textbf{Сводные документы требований к системе}

\textbf{<<Доставка новогодних подарков>>}

\vspace{4mm}

Репозиторий: \url{https://github.com/Zimin0/ITMO-group-12}

\vfill

\end{center}

\begin{flushright}
Выполнили: студенты групп К3420, К3422, К3423\\[2mm]
Зименков Никита К3423\\
Калачева Вера К3423\\
Козлов Никита К3422\\
Луценко Владимир К3420\\
Петрова Виктория К3423\\
Федоров Владислав К3423

\vspace{6mm}

Проверила: доцент\\
Войтюк Татьяна Евгеньевна
\end{flushright}

\vfill

\begin{center}
Санкт-Петербург\\
2026
\end{center}

\end{titlepage}

% ========== СОДЕРЖАНИЕ ==========
\tableofcontents
\newpage

% ====================================================================
%                ЧАСТЬ I. ДОКУМЕНТ О КОНЦЕПЦИИ И ГРАНИЦАХ
% ====================================================================

\section*{Часть I. Документ о концепции и границах}
\addcontentsline{toc}{section}{Часть I. Документ о концепции и границах}

\section{Бизнес-требования}

\subsection{Исходные данные}

\subsubsection{Общее описание проекта}

Проект <<Доставка новогодних подарков>> (Santa CRM) представляет собой клиент-серверное веб-приложение для автоматизации процесса сбора, обработки и доставки новогодних подарков детям. Система охватывает полный цикл от получения письма ребёнка Деду Морозу до вручения подарка курьером.

Проект выполняется в рамках дисциплины <<Создание клиент-серверных приложений>> (7--8 семестры) командой из шести студентов.

Актуальность обусловлена необходимостью координации множества участников процесса доставки подарков: детей, помощников Деда Мороза, складских работников, курьеров и благотворительных организаций.

\subsubsection{Анализ предметной области}

Система ориентирована на автоматизацию благотворительной деятельности в период новогодних праздников. Роли пользователей и их функции описаны в таблице~\ref{tab:roles}.

\begin{table}[H]
\centering
\caption{Роли пользователей системы}
\label{tab:roles}
\small
\begin{tabular}{|L{3.5cm}|L{9.5cm}|}
\hline
\textbf{Роль} & \textbf{Функции} \\
\hline
Ребёнок & Написание писем Деду Морозу, просмотр статуса своих писем \\
\hline
Помощник Деда Мороза & Модерация писем (одобрение/отклонение), просмотр отчётов \\
\hline
Складская служба & Управление запасами подарков, обновление количества на складе \\
\hline
Курьер & Просмотр назначенных доставок, обновление статуса доставки \\
\hline
Благотворительная организация & Просмотр отчётов и статистики \\
\hline
Администратор & Полный доступ ко всем функциям, управление через Django Admin \\
\hline
\end{tabular}
\end{table}

\subsection{Возможности бизнеса}

Система решает следующие задачи:
\begin{itemize}
    \item централизация процесса сбора и обработки детских писем;
    \item автоматизация управления складом подарков;
    \item координация доставки с отслеживанием статусов;
    \item формирование отчётности для благотворительных организаций;
    \item экспорт отчётов в формате PDF.
\end{itemize}

Потенциальные потребители: благотворительные фонды, муниципальные организации, волонтёрские объединения.

\subsection{Бизнес-цели}

\begin{enumerate}
    \item Разработать MVP клиент-серверного приложения, охватывающего полный цикл от получения письма до доставки подарка.
    \item Реализовать ролевую модель доступа с шестью типами пользователей.
    \item Обеспечить прозрачную отчётность по всем этапам процесса.
    \item Обеспечить возможность экспорта данных в PDF для внешних стейкхолдеров.
\end{enumerate}

Целевые показатели MVP: поддержка 6 ролей, обработка писем с 5 статусами жизненного цикла, управление заказами с 6 статусами, логистика с 4 статусами доставки.

\subsection{Критерии успеха}

\begin{itemize}
    \item полнота реализации основного функционала (письма, склад, логистика, отчётность);
    \item корректная работа ролевой модели доступа;
    \item успешная демонстрация всех пользовательских сценариев;
    \item стабильная работа клиент-серверного взаимодействия;
    \item наличие сервисного слоя для бизнес-логики.
\end{itemize}

\subsection{Видение решения}

Система <<Доставка новогодних подарков>> -- веб-приложение с клиент-серверной архитектурой, реализованное на фреймворке Django. Сервер обрабатывает HTTP-запросы, управляет данными через ORM и возвращает HTML-страницы. Клиент взаимодействует с сервером через браузер.

Основной сценарий работы:
\begin{enumerate}
    \item Ребёнок регистрируется и пишет письмо Деду Морозу.
    \item Помощник модерирует письмо (одобряет или отклоняет).
    \item На основе одобренного письма формируется заказ подарков.
    \item Складская служба комплектует заказ.
    \item Курьер получает назначение и доставляет подарок.
    \item Благотворительные организации и администратор отслеживают статистику.
\end{enumerate}

\subsection{Риски проекта}

\begin{table}[H]
\centering
\caption{Риски проекта}
\label{tab:risks}
\small
\begin{tabular}{|L{3.5cm}|C{2cm}|C{2cm}|L{5cm}|}
\hline
\textbf{Риск} & \textbf{Вероят\-ность} & \textbf{Влияние} & \textbf{Меры митигации} \\
\hline
Нехватка времени на реализацию & Средняя & Высокое & Приоритизация MVP, итеративная разработка \\
\hline
Ошибки в ролевой модели доступа & Низкая & Высокое & Тестирование доступа для каждой роли \\
\hline
Потеря данных (SQLite) & Низкая & Среднее & Резервное копирование, миграции Django \\
\hline
Конфликты при слиянии веток & Средняя & Низкое & Код-ревью, ветка dev для интеграции \\
\hline
Некорректная генерация PDF & Низкая & Низкое & Тестирование шаблонов xhtml2pdf \\
\hline
\end{tabular}
\end{table}

\subsection{Предположения и зависимости}

Предположения: пользователи имеют доступ к веб-браузеру; количество одновременных пользователей не превышает нескольких десятков; данные хранятся локально в SQLite.

Зависимости: Python 3.10+, Django 5.2.12, xhtml2pdf, встроенная система аутентификации Django.

\newpage

\section{Рамки и ограничения проекта}

\subsection{Контекстная диаграмма}

На рисунке~\ref{fig:context} представлена контекстная диаграмма системы.

\begin{figure}[H]
\centering
\begin{tikzpicture}[scale=0.95, transform shape,
    actor/.style={rectangle, draw, rounded corners=3pt, minimum width=3cm, minimum height=0.8cm, font=\small},
    system/.style={rectangle, draw, fill=blue!10, minimum width=4cm, minimum height=1.5cm, font=\bfseries},
    db/.style={cylinder, draw, fill=gray!10, shape border rotate=90, aspect=0.2, minimum width=2cm, minimum height=1cm, font=\small},
    arrow/.style={-{Stealth[length=2mm]}, thick},
    every node/.style={font=\small}
]
    % Центральная система
    \node[system] (sys) {Santa CRM};
    \node[db, right=2.5cm of sys] (db) {SQLite};

    % Акторы
    \node[actor, above left=1.5cm and 3cm of sys] (child) {Ребёнок};
    \node[actor, left=3.5cm of sys] (helper) {Помощник};
    \node[actor, below left=1.5cm and 3cm of sys] (warehouse) {Склад};
    \node[actor, above right=1.5cm and 3cm of sys] (courier) {Курьер};
    \node[actor, below right=1.5cm and 3cm of sys] (charity) {Благотв. орг.};
    \node[actor, below=2cm of sys] (admin) {Администратор};

    % Стрелки
    \draw[arrow, <->] (child) -- node[above, sloped, font=\scriptsize] {письма} (sys);
    \draw[arrow, <->] (helper) -- node[above, font=\scriptsize] {модерация} (sys);
    \draw[arrow, <->] (warehouse) -- node[below, sloped, font=\scriptsize] {запасы} (sys);
    \draw[arrow, <->] (courier) -- node[above, sloped, font=\scriptsize] {доставки} (sys);
    \draw[arrow, <->] (charity) -- node[below, sloped, font=\scriptsize] {отчёты} (sys);
    \draw[arrow, <->] (admin) -- node[right, font=\scriptsize] {управление} (sys);
    \draw[arrow, <->] (sys) -- node[above, font=\scriptsize] {данные} (db);
\end{tikzpicture}
\caption{Контекстная диаграмма системы}
\label{fig:context}
\end{figure}

Каждая внешняя сущность взаимодействует с системой через веб-интерфейс (HTTP-запросы к серверу Django). Сервер обрабатывает запросы, обращается к базе данных SQLite и возвращает HTML-ответы.

\subsection{Дерево функций}

Основные функциональные блоки системы:
\begin{enumerate}
    \item \textbf{Аутентификация и авторизация}: регистрация, вход, выход, ролевой доступ.
    \item \textbf{Управление письмами}: создание (ребёнок), просмотр, модерация (помощник/админ), фильтрация.
    \item \textbf{Каталог подарков}: просмотр списка подарков и категорий.
    \item \textbf{Управление складом}: просмотр запасов, обновление количества.
    \item \textbf{Логистика}: просмотр доставок, обновление статуса, фиксация даты вручения.
    \item \textbf{Отчётность}: дашборд со статистикой, экспорт в PDF.
    \item \textbf{Администрирование}: Django Admin для управления всеми сущностями.
\end{enumerate}

\subsection{Список событий системы}

\begin{table}[H]
\centering
\caption{События системы}
\label{tab:events}
\small
\begin{tabular}{|C{0.7cm}|L{3.5cm}|L{8.5cm}|}
\hline
\textbf{ID} & \textbf{Событие} & \textbf{Реакция системы} \\
\hline
E1 & Регистрация & Создание учётной записи с ролью, автоматический вход \\
\hline
E2 & Авторизация & Создание сессии, отображение интерфейса по роли \\
\hline
E3 & Создание письма & Сохранение со статусом <<Новое>> \\
\hline
E4 & Модерация письма & Смена статуса на <<Одобрено>>/<<Отклонено>> \\
\hline
E5 & Просмотр каталога & Отображение списка подарков/категорий \\
\hline
E6 & Обновление запасов & Изменение количества товара на складе \\
\hline
E7 & Обновление доставки & Смена статуса, при <<Доставлено>> фиксация даты \\
\hline
E8 & Просмотр отчётов & Формирование статистики из БД \\
\hline
E9 & Экспорт PDF & Генерация PDF через xhtml2pdf \\
\hline
\end{tabular}
\end{table}

\subsection{Состав выпусков}

\textbf{MVP (текущий выпуск):} регистрация/аутентификация с 6 ролями; письма с модерацией (5 статусов); каталог подарков и категорий; управление складскими запасами; логистика (4 статуса); дашборд с экспортом в PDF; панель администратора.

\textbf{Возможный второй выпуск:} автоматическое формирование заказов; уведомления по email; миграция на PostgreSQL; REST API для мобильного клиента.

\subsection{Ограничения и исключения}

Ограничения MVP: только веб-интерфейс, SQLite (ограничения параллельного доступа), нет интеграции с внешними сервисами, сервер разработки Django, интерфейс только на русском.

Исключения: платёжная система, чат между пользователями, геолокация курьеров, интеграция с внешними складскими системами.

\newpage

\section{Бизнес-контекст}

\subsection{Профили заинтересованных лиц}

\begin{table}[H]
\centering
\caption{Заинтересованные лица}
\label{tab:stakeholders}
\small
\begin{tabular}{|L{3cm}|L{4.5cm}|L{5cm}|}
\hline
\textbf{Стейкхолдер} & \textbf{Интересы} & \textbf{Влияние на проект} \\
\hline
Ребёнок & Удобная отправка письма & Основной пользователь \\
\hline
Помощник & Удобная модерация & Контроль качества данных \\
\hline
Склад. служба & Актуальная информация & Обеспечение наличия подарков \\
\hline
Курьер & Чёткие задания & Вручение подарка \\
\hline
Благотв. орг. & Прозрачная отчётность & Финансирование, контроль \\
\hline
Администратор & Стабильная работа & Технические решения \\
\hline
Преподаватель & Соответствие требованиям & Оценка проекта \\
\hline
\end{tabular}
\end{table}

\subsection{Персона пользователя}

\begin{table}[H]
\centering
\caption{Персона: Ребёнок}
\label{tab:persona_child}
\small
\begin{tabular}{|L{3.5cm}|L{9cm}|}
\hline
Имя & Маша \\
\hline
Возраст & 8 лет \\
\hline
Город & Санкт-Петербург \\
\hline
Цель & Написать письмо Деду Морозу \\
\hline
Ожидания & Простая форма, понятный интерфейс \\
\hline
Опасения & Письмо не дойдёт \\
\hline
\end{tabular}
\end{table}

\begin{table}[H]
\centering
\caption{Персона: Помощник Деда Мороза}
\label{tab:persona_helper}
\small
\begin{tabular}{|L{3.5cm}|L{9cm}|}
\hline
Имя & Алексей \\
\hline
Возраст & 25 лет \\
\hline
Роль & Волонтёр, модерирует письма \\
\hline
Цель & Быстро обрабатывать входящие письма \\
\hline
Ожидания & Фильтрация по статусу \\
\hline
\end{tabular}
\end{table}

\subsection{Приоритеты проекта}

\begin{table}[H]
\centering
\caption{Приоритеты функциональности}
\label{tab:priorities}
\small
\resizebox{\textwidth}{!}{%
\begin{tabular}{|C{1.2cm}|L{5.6cm}|L{5.2cm}|}
\hline
\textbf{Приор.} & \textbf{Функциональность} & \textbf{Обоснование} \\
\hline
P0 & Аутентификация и ролевой доступ & Основа безопасности \\
\hline
P0 & Письма с модерацией & Ядро бизнес-логики \\
\hline
P0 & Управление складом & Обеспечение наличия подарков \\
\hline
P1 & Логистика доставки & Завершение цикла \\
\hline
P1 & Отчётность и PDF-экспорт & Прозрачность для стейкхолдеров \\
\hline
P2 & Каталог и админ-панель & Справочная информация \\
\hline
\end{tabular}}
\end{table}

\newpage

% ====================================================================
%          ЧАСТЬ II. СПЕЦИФИКАЦИЯ ТРЕБОВАНИЙ К ПО (SRS)
% ====================================================================

\section*{Часть II. Спецификация требований к ПО}
\addcontentsline{toc}{section}{Часть II. Спецификация требований к ПО (SRS)}

\section{Общее описание}

\subsection{Операционная среда}

\begin{table}[H]
\centering
\caption{Операционная среда}
\label{tab:env}
\small
\begin{tabular}{|L{4cm}|L{9cm}|}
\hline
\textbf{Компонент} & \textbf{Описание} \\
\hline
Язык & Python 3.10+ \\
\hline
Фреймворк & Django 5.2.12 \\
\hline
СУБД & SQLite 3 \\
\hline
Сервер & Django development server \\
\hline
Клиент & Веб-браузер \\
\hline
ОС & Кроссплатформенная \\
\hline
Библиотеки & xhtml2pdf (PDF-отчёты) \\
\hline
\end{tabular}
\end{table}

\subsection{Выбор средств реализации}

\textbf{Django} выбран благодаря встроенной системе аутентификации с кастомной моделью пользователя, ORM для работы с БД, системе миграций, шаблонному движку и админ-панели.

\textbf{SQLite} выбрана для MVP из-за простоты развёртывания (не требуется отдельный сервер БД).

\textbf{Архитектура} построена по паттерну MVT (Model--View--Template) с дополнительным сервисным слоем (services.py) для бизнес-логики.

\subsection{Пользовательские сценарии}

\subsubsection{UC-1: Регистрация пользователя}

\begin{table}[H]
\centering
\small
\begin{tabular}{|L{3cm}|L{10cm}|}
\hline
Название & Регистрация нового пользователя \\
\hline
Актёр & Незарегистрированный пользователь \\
\hline
Предусловие & Пользователь не авторизован \\
\hline
Основной поток & 1. Открывает /register/. 2. Заполняет форму (логин, пароль, роль, телефон, адрес). 3. Нажимает <<Зарегистрироваться>>. 4. Система создаёт аккаунт и выполняет вход. \\
\hline
Постусловие & Пользователь авторизован \\
\hline
Исключения & Логин уже занят \\
\hline
\end{tabular}
\end{table}

\subsubsection{UC-2: Написание письма Деду Морозу}

\begin{table}[H]
\centering
\small
\begin{tabular}{|L{3cm}|L{10cm}|}
\hline
Актёр & Ребёнок \\
\hline
Предусловие & Авторизован с ролью <<Ребёнок>> \\
\hline
Основной поток & 1. Нажимает <<Написать письмо>>. 2. Заполняет текст. 3. Отправляет. 4. Система сохраняет со статусом <<Новое>>. \\
\hline
Постусловие & Письмо в списке писем ребёнка \\
\hline
\end{tabular}
\end{table}

\subsubsection{UC-3: Модерация письма}

\begin{table}[H]
\centering
\small
\begin{tabular}{|L{3cm}|L{10cm}|}
\hline
Актёр & Помощник / Администратор \\
\hline
Предусловие & Есть письма со статусом <<Новое>> \\
\hline
Основной поток & 1. Открывает список писем. 2. Выбирает письмо. 3. Читает содержание. 4. Одобряет или отклоняет. 5. Система обновляет статус. \\
\hline
Постусловие & Статус изменён, модератор зафиксирован \\
\hline
\end{tabular}
\end{table}

\subsubsection{UC-4: Управление складом}

\begin{table}[H]
\centering
\small
\begin{tabular}{|L{3cm}|L{10cm}|}
\hline
Актёр & Складская служба / Администратор \\
\hline
Основной поток & 1. Открывает страницу склада. 2. Видит товары с количеством. 3. Вводит новое количество. 4. Система сохраняет. \\
\hline
Постусловие & Количество на складе обновлено \\
\hline
\end{tabular}
\end{table}

\subsubsection{UC-5: Обновление статуса доставки}

\begin{table}[H]
\centering
\small
\begin{tabular}{|L{3cm}|L{10cm}|}
\hline
Актёр & Курьер / Администратор \\
\hline
Основной поток & 1. Открывает список доставок. 2. Обновляет статус (<<В пути>>, <<Доставлено>>, <<Не удалось>>). 3. При <<Доставлено>> фиксируется дата. \\
\hline
Постусловие & Статус доставки обновлён \\
\hline
\end{tabular}
\end{table}

\subsubsection{UC-6: Просмотр отчётов и экспорт PDF}

\begin{table}[H]
\centering
\small
\begin{tabular}{|L{3cm}|L{10cm}|}
\hline
Актёр & Администратор / Помощник / Благотв. орг. \\
\hline
Основной поток & 1. Открывает отчёты. 2. Видит дашборд (кол-во писем, заказов, доставок, разбивка по статусам и категориям). 3. Нажимает <<Экспорт PDF>>. 4. Получает файл. \\
\hline
Постусловие & PDF-отчёт скачан \\
\hline
\end{tabular}
\end{table}

\newpage

\section{Функции системы и варианты использования}

\subsection{Функциональные требования}

\begin{table}[H]
\centering
\caption{Функциональные требования}
\label{tab:func_req}
\small
\begin{tabular}{|C{0.8cm}|L{4cm}|L{8cm}|}
\hline
\textbf{ID} & \textbf{Требование} & \textbf{Описание} \\
\hline
FR-1 & Регистрация & Форма с выбором роли (кроме Admin), автовход \\
\hline
FR-2 & Аутентификация & Вход/выход через Django auth \\
\hline
FR-3 & Ролевой доступ & Разграничение по ролям (декоратор role\_required) \\
\hline
FR-4 & Создание письма & Форма для текста, привязка к ребёнку \\
\hline
FR-5 & Модерация писем & Одобрение/отклонение, запись модератора \\
\hline
FR-6 & Фильтрация писем & Фильтр по статусу (GET-параметр) \\
\hline
FR-7 & Каталог подарков & Список с категорией и ценой \\
\hline
FR-8 & Каталог категорий & Список с описанием \\
\hline
FR-9 & Управление запасами & Просмотр/обновление количества \\
\hline
FR-10 & Управление доставками & Список доставок, смена статуса \\
\hline
FR-11 & Дашборд отчётности & Статистика по письмам, заказам, доставкам \\
\hline
FR-12 & Экспорт PDF & Генерация через xhtml2pdf \\
\hline
FR-13 & Админ-панель & Django Admin \\
\hline
\end{tabular}
\end{table}

\subsection{Матрица <<Роль -- Функция>>}

\begin{table}[H]
\centering
\caption{Матрица доступа}
\label{tab:access_matrix}
\small
\begin{tabular}{|L{2.8cm}|C{1.2cm}|C{1.2cm}|C{1.2cm}|C{1.2cm}|C{1.2cm}|C{1.2cm}|}
\hline
\textbf{Функция} & \rotatebox{90}{\small\textbf{Ребёнок}} & \rotatebox{90}{\small\textbf{Помощник}} & \rotatebox{90}{\small\textbf{Склад}} & \rotatebox{90}{\small\textbf{Курьер}} & \rotatebox{90}{\small\textbf{Благотв.}} & \rotatebox{90}{\small\textbf{Админ}} \\
\hline
Регистрация & + & + & + & + & + & -- \\
\hline
Создание письма & + & & & & & \\
\hline
Модерация & & + & & & & + \\
\hline
Склад & & & + & & & + \\
\hline
Доставки & & + & & + & & + \\
\hline
Отчёты & & + & & & + & + \\
\hline
Каталог & + & + & + & + & + & + \\
\hline
Админ-панель & & & & & & + \\
\hline
\end{tabular}
\end{table}

\newpage

\section{Требования к данным}

\subsection{Логическая модель данных}

Система содержит 7 основных сущностей, связанных через Django ORM. ER-диаграмма представлена на рисунке~\ref{fig:erd}.

\begin{figure}[H]
\centering
\begin{tikzpicture}[
    entity/.style={rectangle, draw, fill=blue!8, minimum width=2.2cm, minimum height=0.9cm, font=\small\bfseries},
    rel/.style={font=\scriptsize, fill=white, inner sep=1pt},
    arrow/.style={-{Stealth[length=2mm]}, thick},
    every node/.style={font=\small}
]
    % Сущности - верхний ряд
    \node[entity] (user) at (0,0) {User};
    \node[entity] (letter) at (4,0) {Letter};
    \node[entity] (order) at (8,0) {GiftOrder};
    \node[entity] (delivery) at (12,0) {Delivery};

    % Сущности - нижний ряд
    \node[entity] (category) at (0,-3) {Category};
    \node[entity] (item) at (4,-3) {Item};
    \node[entity] (inventory) at (8,-3) {Inventory};

    % Связи - верхний ряд
    \draw[arrow] (user) -- node[rel, above] {1 : N} node[rel, below] {\scriptsize child} (letter);
    \draw[arrow] (letter) -- node[rel, above] {1 : 1} (order);
    \draw[arrow] (order) -- node[rel, above] {1 : 1} (delivery);

    % User -> Letter (moderated_by) и User -> Delivery (courier)
    \draw[arrow, dashed] (user) .. controls (2, 1.2) and (6, 1.2) .. node[rel, above] {\scriptsize moderated\_by} (letter.north);
    \draw[arrow, dashed] (user) .. controls (4, 2) and (10, 2) .. node[rel, above] {\scriptsize courier} (delivery.north);

    % Связи - нижний ряд
    \draw[arrow] (category) -- node[rel, above] {1 : N} (item);
    \draw[arrow] (item) -- node[rel, above] {1 : 1} (inventory);

    % M:N между Item и GiftOrder
    \draw[arrow, <->] (item) -- node[rel, right] {M : N} (order);
\end{tikzpicture}
\caption{ER-диаграмма (логическая модель данных)}
\label{fig:erd}
\end{figure}

\subsection{Словарь данных}

\begin{table}[H]
\centering
\caption{Модель User}
\label{tab:dict_user}
\small
\begin{tabular}{|L{2.2cm}|L{2.5cm}|L{3cm}|L{4.8cm}|}
\hline
\textbf{Поле} & \textbf{Тип} & \textbf{Ограничения} & \textbf{Описание} \\
\hline
id & integer & PK, auto & Первичный ключ \\
\hline
username & varchar(150) & unique, not null & Имя пользователя \\
\hline
password & varchar(128) & not null & Хэш пароля \\
\hline
role & varchar(20) & choices, default=child & child, helper, courier, warehouse, charity, admin \\
\hline
phone & varchar(20) & blank & Телефон \\
\hline
address & text & blank & Адрес \\
\hline
\end{tabular}
\end{table}

\begin{table}[H]
\centering
\caption{Модель Letter}
\label{tab:dict_letter}
\small
\begin{tabular}{|L{2.5cm}|L{2.5cm}|L{3cm}|L{4.5cm}|}
\hline
\textbf{Поле} & \textbf{Тип} & \textbf{Ограничения} & \textbf{Описание} \\
\hline
id & integer & PK, auto & Первичный ключ \\
\hline
child\_id & integer & FK $\to$ User & Автор письма \\
\hline
content & text & not null & Текст письма \\
\hline
status & varchar(20) & choices & new, approved, rejected, processing, completed \\
\hline
created\_at & datetime & auto\_now\_add & Дата создания \\
\hline
moderated\_by & integer & FK $\to$ User, null & Модератор \\
\hline
\end{tabular}
\end{table}

\begin{table}[H]
\centering
\caption{Модель Item}
\label{tab:dict_item}
\small
\begin{tabular}{|L{2.5cm}|L{2.5cm}|L{3cm}|L{4.5cm}|}
\hline
\textbf{Поле} & \textbf{Тип} & \textbf{Ограничения} & \textbf{Описание} \\
\hline
id & integer & PK, auto & Первичный ключ \\
\hline
title & varchar(200) & not null & Название подарка \\
\hline
description & text & blank & Описание \\
\hline
category\_id & integer & FK $\to$ Category & Категория \\
\hline
price & decimal(10,2) & default=0 & Условная стоимость \\
\hline
\end{tabular}
\end{table}

\begin{table}[H]
\centering
\caption{Модели Inventory, GiftOrder, Delivery}
\label{tab:dict_other}
\small
\resizebox{\textwidth}{!}{%
\begin{tabular}{|L{3cm}|L{2cm}|L{2.4cm}|L{4.5cm}|}
\hline
\textbf{Модель.Поле} & \textbf{Тип} & \textbf{Ограничения} & \textbf{Описание} \\
\hline
Inventory.item & integer & 1:1 $\to$ Item & Товар \\
\hline
Inventory.stock & integer & $\geq 0$ & Кол-во на складе \\
\hline
Inventory.\-reserved & integer & $\geq 0$ & Зарезервировано \\
\hline
GiftOrder.letter & integer & 1:1 $\to$ Letter & Связанное письмо \\
\hline
GiftOrder.items & M2M & $\to$ Item & Подарки в заказе \\
\hline
GiftOrder.status & varchar(20) & choices & 6 статусов \\
\hline
Delivery.order & integer & 1:1 $\to$ GiftOrder & Связанный заказ \\
\hline
Delivery.courier & integer & FK $\to$ User & Курьер \\
\hline
Delivery.status & varchar(20) & choices & 4 статуса \\
\hline
Delivery.route\-\_data & JSON & null, blank & Данные маршрута \\
\hline
Delivery.\-delivered\_at & datetime & null & Дата вручения \\
\hline
\end{tabular}}
\end{table}

\subsection{Отчётность}

Система формирует статистический отчёт: общее количество писем, подарков, категорий, заказов, доставок; разбивку писем по статусам; количество товаров по категориям; дату формирования. Доступно в HTML (дашборд) и PDF (экспорт через xhtml2pdf).

\subsection{Обеспечение целостности данных}

\begin{itemize}
    \item \textbf{Ссылочная целостность}: внешние ключи с CASCADE или SET\_NULL;
    \item \textbf{Уникальность}: username (User), name (Category);
    \item \textbf{Валидация}: Django Forms (серверная);
    \item \textbf{Неотрицательность}: PositiveIntegerField для запасов;
    \item \textbf{Автоматические поля}: auto\_now\_add, auto\_now;
    \item \textbf{Транзакции}: Django ORM оборачивает операции в транзакции.
\end{itemize}

\newpage

\section{Требования к внешним интерфейсам}

\subsection{Пользовательские интерфейсы}

Интерфейс реализован через Django Templates и CSS. Основные экраны:

\begin{enumerate}
    \item \textbf{Главная} (home.html) -- роль пользователя, навигация по ролям, статистика.
    \item \textbf{Регистрация} (register.html) -- логин, пароль, роль, телефон, адрес.
    \item \textbf{Вход} (login.html) -- стандартная форма Django.
    \item \textbf{Письма} -- список с фильтром, форма создания, детали с модерацией.
    \item \textbf{Склад} (inventory.html) -- товары с остатком, форма обновления.
    \item \textbf{Доставки} (delivery\_list.html) -- список с обновлением статуса.
    \item \textbf{Отчёты} (dashboard.html) -- дашборд + кнопка экспорта PDF.
    \item \textbf{Каталог} -- подарки и категории.
\end{enumerate}

Навигация через верхнюю панель: Главная, Объекты, Категории, Admin, Вход/Выход.

\subsection{Программные интерфейсы}

\begin{table}[H]
\centering
\caption{URL-маршруты системы}
\label{tab:urls}
\small
\begin{tabular}{|L{5cm}|C{1.5cm}|L{5.5cm}|}
\hline
\textbf{URL} & \textbf{Метод} & \textbf{Описание} \\
\hline
/ & GET & Главная \\
\hline
/register/ & GET, POST & Регистрация \\
\hline
/login/, /logout/ & GET, POST & Вход/выход \\
\hline
/letters/ & GET & Список писем \\
\hline
/letters/create/ & GET, POST & Создание письма \\
\hline
/letters/<pk>/moderate/ & GET, POST & Модерация \\
\hline
/warehouse/ & GET & Склад \\
\hline
/warehouse/<id>/update/ & POST & Обновление запасов \\
\hline
/logistics/ & GET & Доставки \\
\hline
/logistics/<pk>/status/ & POST & Обновление статуса \\
\hline
/reports/ & GET & Отчёты \\
\hline
/reports/export/ & GET & Экспорт PDF \\
\hline
/items/, /categories/ & GET & Каталог \\
\hline
/admin/ & GET & Админ-панель \\
\hline
\end{tabular}
\end{table}

\subsection{Архитектура взаимодействия клиента и сервера}

\begin{figure}[H]
\centering
\begin{tikzpicture}[
    block/.style={rectangle, draw, fill=blue!8, minimum width=2.8cm, minimum height=0.8cm, font=\small, text width=2.6cm, align=center},
    arrow/.style={-{Stealth[length=2mm]}, thick},
    every node/.style={font=\small}
]
    % Клиент
    \node[block, fill=green!10] (browser) at (0,0) {Браузер\\(клиент)};

    % Сервер
    \node[block] (urlconf) at (4.5,0) {URLconf\\(urls.py)};
    \node[block] (view) at (9,0) {View\\(views.py)};
    \node[block, fill=orange!10] (auth) at (9,1.8) {Декораторы\\доступа};
    \node[block, fill=yellow!10] (service) at (9,-1.8) {Service Layer\\(services.py)};
    \node[block] (orm) at (13.5,-1.8) {Django ORM};
    \node[block, fill=gray!15] (db) at (13.5,-3.6) {SQLite};
    \node[block, fill=red!8] (template) at (13.5,0) {Template\\Engine};

    % Стрелки
    \draw[arrow, <->] (browser) -- node[above, font=\scriptsize] {HTTP} (urlconf);
    \draw[arrow] (urlconf) -- (view);
    \draw[arrow, <->] (view) -- (auth);
    \draw[arrow, <->] (view) -- (service);
    \draw[arrow, <->] (service) -- (orm);
    \draw[arrow, <->] (orm) -- (db);
    \draw[arrow] (view) -- (template);
    \draw[arrow] (template) |- node[near end, above, font=\scriptsize] {HTML} (browser);
\end{tikzpicture}
\caption{Архитектура клиент-серверного взаимодействия}
\label{fig:architecture}
\end{figure}

\newpage

\section{Атрибуты качества}

\subsection{Внешние атрибуты качества}

\begin{table}[H]
\centering
\caption{Внешние атрибуты качества}
\label{tab:ext_quality}
\small
\begin{tabular}{|L{3.8cm}|L{9cm}|}
\hline
\textbf{Атрибут} & \textbf{Описание и реализация} \\
\hline
Удобство использования (Usability) & Интерфейс адаптируется к роли пользователя: главная страница показывает только релевантные действия. Django messages информируют о результатах. \\
\hline
Функциональная пригодность & Полный цикл обработки: письмо, модерация, заказ, склад, доставка, отчёт. Все 6 ролей имеют свой функционал. \\
\hline
Надёжность (Reliability) & Серверная валидация форм. Ролевой доступ. Внешние ключи с каскадным удалением обеспечивают целостность. \\
\hline
Производительность (Performance) & Использование select\_related в сервисном слое для минимизации SQL-запросов. Лёгкий стек для учебного масштаба. \\
\hline
Безопасность (Security) & Хэширование паролей, CSRF-защита, ролевой доступ, роль Admin исключена из регистрации. \\
\hline
\end{tabular}
\end{table}

\subsection{Внутренние атрибуты качества}

\begin{table}[H]
\centering
\caption{Внутренние атрибуты качества}
\label{tab:int_quality}
\small
\begin{tabular}{|L{3.8cm}|L{9cm}|}
\hline
\textbf{Атрибут} & \textbf{Описание и реализация} \\
\hline
Сопровождаемость (Maintainability) & Разделение ответственности: models.py, views.py, services.py, forms.py, templates/. Единообразное именование. \\
\hline
Переносимость (Portability) & Кроссплатформенный стек (Python + Django). Зависимости зафиксированы в requirements.txt. \\
\hline
Тестируемость (Testability) & Сервисный слой изолирует бизнес-логику от представления. Контрольный список тестов в TEST\_CHECKLIST.md. \\
\hline
Расширяемость (Extensibility) & Модульная архитектура Django. Миграции поддерживают эволюцию схемы БД. \\
\hline
\end{tabular}
\end{table}

\subsection{Интернационализация}

В текущей версии система поддерживает только русский язык. Архитектура допускает i18n: Django имеет встроенную поддержку через gettext, все строки вынесены в шаблоны и verbose\_name моделей.

\newpage

\section*{Заключение}
\addcontentsline{toc}{section}{Заключение}

В рамках практической работы разработан итоговый отчёт, включающий документ о концепции и границах системы <<Доставка новогодних подарков>> и спецификацию требований к программному обеспечению.

Система реализована как клиент-серверное веб-приложение на Django 5.2.12. Реализован полный цикл обработки от создания писем до доставки подарков. Поддерживаются шесть ролей пользователей. Архитектура следует паттерну MVT с сервисным слоем. База данных содержит 7 сущностей. PDF-отчёты генерируются через xhtml2pdf.

Исходный код: \url{https://github.com/Zimin0/ITMO-group-12}.

Видеодемонстрация: \url{https://drive.google.com/drive/folders/1TdcRUdQAou28KZtcZnhB0-76Mlh_jUlx?usp=sharing}.

\newpage

\section*{Список литературы}
\addcontentsline{toc}{section}{Список литературы}

\begin{enumerate}
    \item Вигерс К., Битти Дж. Разработка требований к программному обеспечению. 3-е изд. М.: Русская редакция, 2014. 736~с.
    \item Django Documentation. URL: \nolinkurl{https://docs.djangoproject.com/en/5.2/} (дата обращения: 15.03.2026).
    \item IEEE Std 830-1998. Recommended Practice for Software Requirements Specifications.
    \item xhtml2pdf Documentation. URL: \nolinkurl{https://xhtml2pdf.readthedocs.io/} (дата обращения: 15.03.2026).
\end{enumerate}

\end{document}
